\clearpage
\begin{landscape}
\begingroup
\fontsize{10}{13}\selectfont
\fancyhead[L]{\small NACHLASS. STÖRUNGEN DER PALLAS}
\renewcommand{\headrulewidth}{0pt}
\pdfbookmark[2]{Allgemeine Störungen, Erste Rechnung, [11]}{pallas-V-fragment11}
\begin{center}[11.]\end{center}
\begingroup
\newcommand{\VElevenLongitude}{\mathord{\tikz[x=1ex,y=1ex,baseline=0.08ex,line width=.075ex,line cap=round,line join=round]{\draw (.30,1.05) .. controls (.24,.92) and (.02,.76) .. (.04,.62) .. controls (.07,.43) and (.52,.43) .. (.56,.62) .. controls (.58,.76) and (.37,.94) .. (.30,1.05);\draw (.30,.10)--(.30,.67);\draw (.04,.30)--(.56,.30);}}}
\par\noindent\hspace*{\parindent}[Da die mittlere Länge durch die Relation
\[
\VElevenLongitude=\int n\,\mathrm{d}t+\varepsilon
\]
definirt ist, so zerfallen die Störungen der Epoche in zwei Theile, deren ersterer aus \(\int n\,\mathrm{d}t\), letzterer aus \(\varepsilon\) selbst entspringt. Der erstere ergibt sich aus dem vorigen Artikel, indem man die Störungen des Logarithmen der mittlern Bewegung
\[
\frac{\delta n}{n}=\delta\log\operatorname{hyp}n
\]
nochmals integrirt. Für unser Beispiel entnehmen wir dem vorigen Artikel den zum Argument \(5M'-2M\) gehörigen Coefficienten von \(-\delta\log\operatorname{hyp}n\), der dort mit \(*\) bezeichnet ist, worauf sich ergibt:
\endgroup
\begingroup
\newcommand{\VElevenPoint}{\mathord{.}\mkern5mu}
\newcommand{\VElevenFive}{\VElevenPoint\VElevenPoint\VElevenPoint\VElevenPoint\VElevenPoint}
\begin{center}
\begin{tabular}{@{}l@{\quad}l@{}}
\(*\;\VElevenFive\VElevenFive\VElevenPoint\VElevenPoint\VElevenPoint2{,}23410_n\)&\\
\(\log\left(5\frac{n'}{n}-2\right)\;\VElevenPoint\VElevenPoint8{,}74966_n\)&also das Störungsglied:\\
\cline{1-1}
\multicolumn{1}{r}{\(3{,}48444\)}&\(-3051''{,}00\sin(5M'-2M+69^0\,57'\,58'').\)
\end{tabular}
\end{center}
\endgroup
\par\noindent\hspace*{\parindent}Es ergibt sich hieraus für die Störungen des ersten Theils der Epoche, wobei die Coefficienten der Argumente \(5M'-2M\), \(8M'-3M\), \(10M'-4M\) später mit verbessertem Divisor neu berechnet sind:
\begingroup
\renewcommand{\arraystretch}{1.15}
\begin{longtable}{@{}r@{}l@{}}
\(\int\delta n\,\mathrm{d}t={}\)&\(-11''{,}47\sin(M+93^0\,28'\,27'')\)\\
&\(-6{,}80\sin(2M+244\;6\;47)\)\\
&\(-1{,}07\sin(3M+286\;9\;45)\)\\[3pt]
&\(-0{,}80\sin(M'-4M+155^0\,24'\,31'')\)\\
&\(-2{,}70\sin(M'-3M+190\;5\;33)\)\\
&\(-6{,}44\sin(M'-2M+0\;16\;14)\)\\
&\(-7{,}51\sin(M'-M+127\;15\;14)\)\\
&\(-5{,}25\sin(M'+M+131\;53\;19)\)\\
&\(-1{,}11\sin(M'+2M+248\;38\;11)\)\\
&\(-0{,}38\sin(M'+3M+3\;0\;46)\)\\[3pt]
&\(-0{,}51\sin(2M'-5M+29^0\,29'\,10'')\)\\
&\(-0{,}83\sin(2M'-4M+16\;17\;10)\)\\
&\(-8{,}58\sin(2M'-3M+301\;0\;25)\)\\
&\(-60''{,}23\sin(2M'-2M+325^0\,15'\,27'')\)\\
&\(-731{,}50\sin(2M'-M+135\;51\;55)\)\\
&\(-1{,}69\sin(2M'+M+180\;26\;27)\)\\
&\(-0{,}63\sin(2M'+2M+273\;44\;35)\)\\[3pt]
&\(-0{,}25\sin(3M'-6M+94\;34\;43)\)\\
&\(-0{,}38\sin(3M'-5M+230\;41\;7)\)\\
&\(-3{,}54\sin(3M'-4M+194\;14\;15)\)\\
&\(-12{,}21\sin(3M'-3M+231\;12\;2)\)\\
&\(-67{,}23\sin(3M'-2M+13\;53\;20)\)\\
&\(-568{,}14\sin(3M'-M+145\;6\;27)\)\\
&\(-0{,}56\sin(3M'+M+190\;43\;15)\)\\[3pt]
&\(-1{,}34\sin(4M'-5M+88\;4\;41)\)\\
&\(-2{,}83\sin(4M'-4M+145\;41\;2)\)\\
&\SourceLeaf{531}\(-16''{,}03\sin(4M'-3M+272^0\,42'\,10'')\)\\
&\(-106{,}81\sin(4M'-2M+42\;6\;35)\)\\
&\(-21{,}21\sin(4M'-M+168\;53\;16)\)\\[3pt]
&\(-0{,}48\sin(5M'-6M+344\;7\;4)\)\\
&\(-0{,}96\sin(5M'-5M+76\;38\;4)\)\\
&\(-4{,}67\sin(5M'-4M+175\;2\;21)\)\\
&\(-16{,}34\sin(5M'-3M+303\;38\;18)\)\\
&\(-3121{,}14\sin(5M'-2M+69\;57\;58)\)\\
&\(-3{,}20\sin(5M'-M+192\;27\;56)\)\\[3pt]
&\(-0{,}42\sin(6M'-6M+357\;25\;14)\)\\
&\(-1{,}51\sin(6M'-5M+78\;28\;31)\)\\
&\(-4{,}96\sin(6M'-4M+204\;27\;13)\)\\
&\(-21{,}26\sin(6M'-3M+332\;22\;56)\)\\
&\(-39{,}24\sin(6M'-2M+86\;31\;33)\)\\
&\(-1{,}00\sin(6M'-M+203\;22\;3)\)\\[3pt]
&\(-0{,}52\sin(7M'-6M+355\;56\;19)\)\\
&\(-1{,}63\sin(7M'-5M+110\;12\;38)\)\\
&\(-4{,}75\sin(7M'-4M+233^0\,40'\,55'')\)\\
&\(-58{,}94\sin(7M'-3M+358\;44\;52)\)\\
&\(-3{,}60\sin(7M'-2M+126\;53\;40)\)\\[3pt]
&\(-0{,}60\sin(8M'-6M+16\;54\;19)\)\\
&\(-1{,}56\sin(8M'-5M+138\;29\;50)\)\\
&\(-5{,}35\sin(8M'-4M+263\;25\;0)\)\\
&\(-193{,}69\sin(8M'-3M+27\;39\;28)\)\\[3pt]
&\(-0{,}58\sin(9M'-6M+41\;12\;9)\)\\
&\(-1{,}46\sin(9M'-5M+165\;36\;3)\)\\
&\(-8{,}76\sin(9M'-4M+291\;12\;39)\)\\
&\(-4{,}38\sin(9M'-3M+52\;59\;0)\)\\[3pt]
&\(-0{,}51\sin(10M'-6M+68^0\,51'\,43'')\)\\
&\(-1{,}45\sin(10M'-5M+191\;55\;30)\)\\
&\(-89{,}91\sin(10M'-4M+313\;17\;2)\)\\[3pt]
&\(-1{,}86\sin(11M'-5M+216\;56\;27)\)\\
&\(-7{,}61\sin(11M'-4M+340\;2\;2).]\)
\end{longtable}
\endgroup
\clearpage
\endgroup
\end{landscape}
